This work presents the development of an automated solution for the drafting process of meeting minutes within a legal context, utilizing Large Language Models (LLMs) executed locally. The methodology begins with the creation of a dataset from meeting records provided by the Regional Electoral Court of Ceará (TRE-CE). This is followed by the implementation of a processing (inference) pipeline based on the Map-Reduce paradigm for data synthesis and structuring into JSON format.

The results demonstrate that models such as LLaMA 3 and Gemma 2 stand out in computational efficiency and structural adherence, respectively, when tasked with data synthesis. Meanwhile, models like Mistral v0.3 excel when charged with the structuring task, following a pre-established template.

However, qualitative evaluation conducted via the LLM-as-a-Judge methodology revealed the occurrence of systemic hallucinations in smaller-scale models. This evidence suggests that, while structural automation is feasible, factual integrity still requires additional validation mechanisms to ensure the legal security of the generated documents.

\keywords{LLM. IA. MCP. transformer}